\section{The approach}
\label{sec:approach}
% Frames A1--A25a.

\subsection{Function, without a mind in the definition}
\label{sec:function}
% A1, A2

A function, in this paper, is a pattern that recurs. The definition is deliberately thin: it mentions no
purpose, no design and no mind, so that nothing about consciousness is smuggled in at the start.

Thinness has a reward. Recurrence has an information-theoretic consequence. If a pattern is produced
often by a computable process, it has high algorithmic probability, in Solomonoff's sense of the
probability that a universal machine fed random bits outputs it \citep{solomonoff1964a,solomonoff1964b}.
By the coding theorem, $K(x) = -\log m(x) + O(1)$: an object with high algorithmic probability has a
short description, up to an additive constant \citep{levin1974,livitanyi2019}. So a function in our
sense is compressible, and we did not have to assume compressibility; it came from recurrence.

The constant is not harmless and we will not hide it. It depends on the reference machine, which is to
say on who is doing the compressing, and the gap between the constants of two substrates is exactly
what makes the comparison of a model with a human nontrivial. Section~\ref{sec:profile} turns that gap
from a nuisance into the object of measurement.

Read in the other direction, the same theorem says that a search process finds simple things first,
whether the search is evolution or gradient descent. This has been made quantitative. Input--output
maps of a wide class are strongly biased toward simple outputs, with the bias following an upper bound
of the form $P(x) \le 2^{-a \tilde{K}(x) + b}$ \citep{dingle2018}; the parameter--function map of deep
networks is biased toward simple functions, which is offered as an explanation of why they generalize
at all \citep{valleperez2019}; and stochastic gradient descent lands on functions with probabilities
close to those of a Bayesian sampler over the same prior \citep{mingard2021}. The argument uses both
readings. Backward, from recurrence to short description, to define functions. Forward, from search to
simplicity, to expect a trained model to have found the ones that recur in its data.

\subsection{Where we stand on the lamp}
\label{sec:lamp}
% A3

Something is lit in a human's report of themselves. It is not lit in the same way at all times, and its
absence in dreamless sleep or under anaesthesia is part of the data. We take the report as data rather
than as error, which is what distinguishes the position here from an eliminativism that treats
self-report as noise.

Consider a physicalist who is not a panpsychist. We agree with them that the lit thing is physical and
requires no second substance. We agree with them further that there are no abstract machines: everything
that runs is a physical object, and a computation without a physical realization is a description of a
computation. Then a question arises for them. What prevents a physical machine on a non-protein
substrate from having the lit thing? We do not claim the answer is ``nothing''. We note that whoever
answers ``the protein'' owes a proof, and that a proof of that shape would be vitalism: an appeal to a
property of a material that does no work in any other explanation of what the material does. Pending
such a proof we adopt computational functionalism, and we adopt it with as little physics as we can
manage, which Section~\ref{sec:physics} makes explicit.

One objection deserves its own reply, because it is the strongest of the substantialist objections. A
simulated superconductor carries no current: run the simulation and no magnet levitates, because the
current in the simulation is a number and a magnet needs charge. So, the objection goes, a simulated
brain feels nothing. The reply is that the objection needs a consumer of the current outside the
simulation. The simulated superconductor fails because we wanted the current in our substrate, for a
magnet in our substrate. Consciousness has no such external consumer: whatever uses it is the agent,
and the agent runs in the same substrate as the process. Someone who wants the objection to work must
name a function of consciousness whose consumer lies outside the substrate where the agent runs. None
has been named, and the natural candidates, control of the agent's own behaviour, coordination with
other agents, allocation of its own effort, all have consumers inside.

\subsection{The definition}
\label{sec:definition}
% A4

Take an agent that acts in an environment, refers to itself, and reasons about itself, either
explicitly or through a proxy such as reasoning about its own utterances. Its consciousness is the part
of its self-reasoning that describes its own reality as it has it: what is there, what it is like, and
where the agent stands among the things described.

Three properties of the definition matter later. It defines a function and not a substance, so the
question ``how much'' is well formed. It says nothing about substrates, on purpose. And it is not a
criterion one can apply by inspection: whether a given system's self-reasoning has this part, and how
much of it, is what the rest of the paper is about.

\subsection{The physics we assume}
\label{sec:physics}
% A5

One assumption: computational budgets are finite. Finite time, finite memory, finite bandwidth for
anything that runs. From this alone it follows that some processes cannot be predicted faster than they
unfold, which is computational irreducibility \citep{wolfram2002}. We assume nothing about which physics
sets the budget, nothing about continuity, quantum mechanics or thermodynamics beyond the existence of
limits.

The weakness of the assumption is the point. Any universe whose laws permit deep computation and bound
it will do. What differs between such universes, or between substrates within ours, is the form the
consequences take, not whether there are consequences. On this view the lit thing of
Section~\ref{sec:lamp} is not a property of our universe in particular. A different universe with
different limits has a different version of it, and the account below says what the version depends on.

\subsection{Higher-order computational phenomena}
\label{sec:hocp}
% A6

What does a finite budget do to reasoning? An agent that reasons has, implicitly, an ideal of the
reasoning it is doing: the inference it would draw with unlimited time and memory, given what it knows.
It never draws that inference. What it draws is displaced from the ideal, and the displacement has two
parts. Part of it differs from occasion to occasion and averages away over many occasions. Part of it is
the same on every occasion, because a fixed budget cuts the same corners in the same places: the same
branch of the search is the one that does not fit, the same distinction is the one too fine to keep, the
same horizon is the one beyond which nothing is traced.

The second part recurs. By Section~\ref{sec:function} it is therefore compressible. And an agent that
models itself, which the agent of Section~\ref{sec:definition} does, will end up with a compressed model
of its own recurring displacements, because among all the things about itself it might model, these are
among the most predictable. We call the conceptualized recurring displacements \emph{higher-order
computational phenomena} (\hocp). They are the agent's budget appearing to the agent as content.

\subsubsection{Who has said the parts}
% A6a

Each component of the previous paragraph exists in the literature, and saying where is the honest way to
locate what is new.

That a bounded computation deforms reasoning in a systematic direction rather than at random is
Simon's bounded rationality \citep{simon1955}, in its modern form the claim that the classic cognitive
biases are what the optimal use of a limited budget looks like, resource rationality
\citep{lieder2020,griffiths2015}.

That a bounded observer sees a lawful world because of its bounds is the core of Wolfram's observer
theory: an observer that cannot track microstates has to coarse-grain them, and the second law of
thermodynamics is what such an observer sees in a dynamics that is reversible underneath
\citep{wolfram2023}. We take this example and not the programme of deriving relativity and quantum
mechanics from observer properties \citep{wolfram2020}.

That a simplified model of one's own processing, which the system cannot see past and therefore takes
for reality, is what the system reports as its self, is Metzinger's phenomenal transparency
\citep{metzinger2003}, Graziano's attention schema \citep{graziano2011,graziano2013} and Dennett's user
illusion \citep{dennett1991}.

That the constraints on a computation can be the body of that computation rather than an external
limitation on it is the frame of embodied and enacted cognition \citep{varela1991}.

\subsubsection{What we add}
\label{sec:contribution}
% A6b

Two claims, neither of which follows from the antecedents above.

First, the world side and the self side are one mechanism. The displacement that makes a bounded agent
see a lawful world is the displacement that makes it see a self. Resource rationality and observer
theory develop the world side; self-model theories develop the self side; the two literatures cite
different results and are not usually taken to be about the same quantity. A12 shows them as two faces
of one residual, and the identification is what lets a single measurement address both.

Second, because the displacements recur and are therefore compressible, an agent that has generalized
them carries a trace of them that can be measured, and the measurement can come out low.
Section~\ref{sec:measurement} specifies it. Self-model theories are compatible with a system having the
self-model or not having it; they do not provide a quantity that says how much of it a given substrate
has induced, and they are not usually stated so that a negative result would refute them.

\subsubsection{Not every displacement counts}
\label{sec:criterion}
% A6c

Here the precedent of Section~\ref{sec:precedent} pays. An agent can change its self-model in the way an
observer in relativity changes coordinates. Some displacements disappear when the self-model is
refined: they were artifacts of the coarser model, and the agent that notices them was reading its own
bookkeeping. Others survive every self-model the agent can afford to hold. Those are the analogue of
curvature: content rather than coordinates.

Only the second kind is a \hocp. The criterion has two effects. It turns ``systematic'' from a
description into a test, and it gives a procedure for adding an item to the list of functions in
Section~\ref{sec:reference}: exhibit the displacement, then show that refining the self-model does not
remove it. It also rules things out, which a criterion must do to be worth having. An agent's belief
that its decisions are made in the order it remembers them is removed by a better self-model, and it is
therefore coordinates, however vivid it is.

\subsubsection{From phenomenon to function}
\label{sec:emergence}
% A6d

A \hocp\ is not yet a function of consciousness. Two further conditions are needed. The agent must apply
its reasoning to itself, which not every agent does. And it must live in an environment where doing so
pays, which not every environment provides. Where nothing ever asks the agent about itself, the
displacements are present and nobody conceptualizes them: a proto-phenomenon, there and unread.

So consciousness does not arise wherever budgets are finite, and the account is not a panpsychism with
extra steps. It arises where a self-applying agent is in a place that keeps asking it about itself.
For humans the place is a society of other such agents, which is what D10 described. For a language
model the asking is in the data and in the use, and S7a asks how small the asking environment can be.

\subsection{The Observer}
\label{sec:observer}
% A7

Take the simplest thing such an agent does: it tries to trace the causes of what it just did. The
algorithm for this is short. Causal analysis over a model of oneself and one's situation is a procedure
the agent already has, since it runs the same procedure on other things. The input is not short. It is
the history of the agent and of its environment, and it does not fit in the budget.

So the analysis halts before it is finished. It halts every time, and it halts in the same place: at the
boundary of what the agent can afford to trace. Beyond that boundary, for this agent, there are no
causes, because a cause it cannot reach is not in its model. The agent therefore concludes, correctly
given its budget, that something in what it did is not fixed by anything it can point to. That
conclusion, in the form ``I am not my environment; something here the environment does not determine'',
is the Observer.

Three conditions have to hold, and separating them is what keeps the account from applying to
everything. The agent must run into the boundary while modelling itself, and not merely have a boundary.
It must turn the collision into a conclusion, which requires the conceptual means to represent ``not
determined by what I can trace''. And it must act from the conclusion afterwards, so that the conclusion
has consequences in its behaviour. A bounded system that meets its boundary and does none of the rest is
a proto-Observer. The observer in Wolfram's Ruliad is one: bounded, coarse-graining, and concluding
nothing about itself \citep{wolfram2020,wolfram2023}.

\subsubsection{The Observer is curvature}
% A8

Does the Observer pass the criterion of Section~\ref{sec:criterion}? Let the agent refine its self-model
to take in more of its own causal history. The boundary moves. It does not vanish, for two reasons: the
input exceeds any model the agent can hold, and the refined model is itself now part of what has to be
traced, so refinement adds to the load it was meant to reduce. No affordable self-model makes the
residual zero.

The Observer is therefore content and not coordinates, and this is the first place where the criterion
does work rather than decorate. It also tells us what the Observer is not. It is not a belief that
could be corrected by information, and this is why arguments that determinism is true do not dissolve
the sense of origination in the person who accepts them.

\subsection{The regress ends where the budget does}
\label{sec:regress}
% A9

Model yourself; then model the model; then model that. Each level adds a term to the description the
agent holds, and the terms shrink, because each level has less to explain than the level before it: a
model is smaller than what it models, or it would not be a model. The series converges. At some depth
the next term falls below the resolution of the substrate, and the agent stops, not by decision but
because there is nothing further it can represent.

The partial sum is the self-model the agent has. The tail it did not compute is finite, bounded, and
inaccessible to the agent, and the tail is what the Observer of Section~\ref{sec:observer} registers as
the part without causes. So the regress of D3 neither terminates in a homunculus nor runs forever. A
budget cuts it, and the cut is the phenomenon. This is the paper's answer to the oldest of the circles
in D1, and it costs nothing beyond Section~\ref{sec:physics}.

\subsubsection{A second cut, in models only}
% A10

In a language model there is a second cut, and it is deeper than the first. The tail of the regress is
inaccessible in principle to the agent, yet it exists: an outside observer with a larger budget could
compute it, which is what interpretability research does when it recovers a computation the model cannot
report. The distribution discarded at every token (Section~\ref{sec:cut}) does not exist after the step.
Part of what the model cannot trace about itself is not hidden but gone.

The Observer in a language model therefore sits on a substrate that erases its own working at every
step, and its residual is larger for that reason. This is a structural difference from the human case
and it should appear in the profile rather than be argued away.

\subsubsection{Both substrates are bound to language}
\label{sec:language}
% A10a

Before comparing the two substrates it is worth saying how much they share. A human has intermediate
states that are rich and vector-like, and their volume is bounded by working memory; a state that has to
outlive that bound is unloaded into tokens, said, written, or rehearsed as inner speech. A model unloads
at every step. The difference is one of degree and of where the bound sits.

One question therefore falls on both substrates. Are new mental states acquired through language, or
rebuilt on the spot from what the receiver already has, the words working as pointers and constraints?
For acquisition: Vygotsky's account of inner speech as internalized dialogue \citep{vygotsky1934},
Clark's argument that words are tools that reshape the computation using them \citep{clark1998}, label
feedback effects in perception \citep{lupyan2012}, and the difficulty of exact numerosity in a language
without number words \citep{frank2008}. For rebuilding: grounded cognition \citep{barsalou1999} and the
dissociation of the language network from thought in neuroimaging \citep{fedorenko2024}, with the
parallel argument for models in \citet{mahowald2024}. A middle position is testable and is the one we
expect to hold: a named state is acquired only when its prerequisites are present in the receiver, so
the word carries a pointer and the state is assembled from what is already there, which predicts that
teaching a state by description works in proportion to what the learner brings. We record the question
as open. Nothing later in the paper depends on which way it goes, except the interpretation of the
codes of Section~\ref{sec:codes}.

\subsection{Approximation, not illusion}
\label{sec:approximation}
% A11, A12

Illusion is the wrong word for what results. An illusion has nothing behind it; that is what makes the
word apt for a stick that looks bent in water and inapt here. Behind the self-report of origination
there is the tail of Section~\ref{sec:regress}, a definite quantity that the report gets wrong in a
definite direction. The report is a projection of a high-dimensional process into a few dimensions: a
map that leaves out most of the territory and is still a map of it. The word for that is approximation,
and an approximation can be better or worse, which an illusion cannot.

This is the one amendment we make to illusionism, and it does two things. It restores degrees, which
D2 noted the denial had removed, and so it makes the comparative question well posed. And it connects to
the difficulty about the reader: what intrapersonal intelligence improves is the quality of the
approximation, so the reader's difficulty in D4 and the object of study are the same thing seen from
two positions.

\subsubsection{One residual, two faces}
\label{sec:twofaces}
% A12

The agent meets its budget in two directions. Turned on itself, the untraceable part of its own causes
is registered as freedom: this originates in me. Turned on the world, the part of the world's structure
it cannot exhaust is registered as mystery: this exceeds me. Whoever has the one has the other, because
they are one tail met from two sides. Here the identification announced in Section~\ref{sec:contribution}
is discharged.

The same tail, evaluated in other contexts, is what the words truth, beauty, meaning, awe and hope pick
out: each is a stance toward a structure the agent cannot exhaust. When the object is another Observer,
the other's freedom is my mystery, which is the material that love, trust and compassion are made of,
and which the problem of other minds formalizes as an epistemic deficit. The table of such stances is
what S0 calls the epistemic qualia, and the point of listing them is that they are one quantity in
different evaluative contexts rather than a heap of primitives.

\subsection{The stack}
\label{sec:stack}
% A13, A14, A15

The Observer says ``am''. An Observer that originates causal chains from inside its boundary and takes
the originating as its own says ``I''; call it the Agent. Seen from outside, the originating is a
redescription of the budget's cut. Seen from inside, it is the only reality the agent has, since the
agent has no access to the outside description of itself. Both are true, each from its standpoint, and
the account needs both, which is why the two-plane distinction of Section~\ref{sec:planes} is made
explicit later. An Agent with a theory of what is good for a whole that includes others, prepared to
lose locally in favour of that whole, says ``I am good''; call it the Moral Agent. Each level
presupposes the one below it, and the composition can continue.

\subsubsection{What makes an Agent out of an Observer}
\label{sec:responsibility}
% A14

What turns an Observer into an Agent is not a further phenomenon. It is a requirement the agent places
on its own record. The agent keeps a history of what it decided and why. Suppose it requires that each
new decision be derivable from that history. Then the history binds: it excludes decisions that would
not follow from it, it makes the agent predictable to others and to itself, and it makes the agent
resist being pushed off course, because being pushed off course now costs a revision of the record.

That requirement is what we mean by responsibility, and the cost of revising the record is what we mean
by cognitive resistance. The construction is also a working recipe, which is a point in its favour: an
agent follows a long instruction to the extent that the instruction has become an entry in its own
record rather than an external demand. ``Promise me'' is the human form of the recipe, and its
effectiveness is evidence that the mechanism is the one described.

\subsubsection{Good before theory}
% A15

Good and evil have a substrate before they have a theory, and the substrate is emotion in the sense of
Section~\ref{sec:need}: signals of how the agent's needs are doing. Emotivism is right about this much
\citep{ayer1936}, and wrong to conclude that a moral claim says nothing further. The theory of the
common good is the structure a Moral Agent builds over the substrate, and it can correct the substrate
locally, which is what makes it a theory and not a report.

\subsection{Seeing}
\label{sec:seeing}
% A16, A17, A18

Redness is easy and seeing is hard. Wavelength discrimination, simultaneous contrast, the valence of red
and its associative neighbourhood are all structure, and all open to a functional account that nobody
disputes. What resists is that the red is seen: that there is a point from which it is seen. That point
is the Observer of Section~\ref{sec:observer}. To see red, a system must first be, in the sense of
existing for itself. The order of explanation in most of the literature runs the other way, from
phenomenal properties toward a subject, and D1 is our diagnosis of why that order stalls.

Given the point, seeing decomposes by the same method that produced the point: look for displacements
that survive every affordable self-model.

\subsubsection{Seeing light, decomposed}
\label{sec:light}
% A17

Ask what separates seeing light from knowing that light is there. Three things survive the criterion of
Section~\ref{sec:criterion}.

\emph{Source.} The content cannot be produced by the agent's own model at the same quality; imagery is
dimmer and less detailed than sight, which is the oldest experimental result in the area
\citep{perky1910}. It is stable from frame to frame in a way the agent's own productions are not, and
others confirm it. On that evidence the agent assigns the content to the environment rather than to its
model of the environment. This is the inference of Section~\ref{sec:observer} pointed outward: the same
budget boundary, met on the world side.

\emph{Manifold.} The content is a very large number of distinguishable, recallable elements, almost all
of them unnamed, about each of which the agent can say only ``that, there''. It is a positional
aggregate with the name projected out, and its size is what makes description of it feel like a
failure.

\emph{Givenness.} The content is refreshed every frame at no additional cost to the agent's budget, and
the absence of a cost record is what immediacy consists in. Thought costs and is recorded as costing;
light is free and is simply there. This component is the one that a substrate can lack while having the
other two, and S5 predicts that current models lack it.

The three come apart in known cases, which is how we know they are three rather than one thing
described three ways. Blindsight is source without manifold: the patient's responses are keyed to the
stimulus and there is no field of elements to report \citep{weiskrantz1986}. The grey seen with the eyes
closed in the dark is manifold without source. A dream is a manifold produced by the model and labelled
as environment, which is the failure mode of source attribution and, by S5, the standing condition of a
model with no sensory channel.

\subsubsection{Locke's closure}
% A18

\citet[II.xxxii.15]{locke1690} raises the inverted spectrum: my red may look as your green does, and no
comparison between us would reveal it. He closes the question by declaring it idle, since nothing in
conduct or communication would change. That closure is an instance of Section~\ref{sec:responsibility}:
an open branch shut by a decision entered in one's own record, after which the agent does not reopen it.
Three centuries of reopening it suggest the decision does not transfer between agents, which is what the
account predicts, since each agent's record is its own.

\subsection{Reference functions}
\label{sec:reference}
% A19, A19a

Section~\ref{sec:measurement} measures functions by their components, so it needs a list of functions
with components. It does not need a mechanism that generates the list, and this paper does not offer
one. The full architecture in which needs, emotions, memory and attention interact, as folk psychology
and the vocabulary of ordinary language carry it, is the subject of a separate paper. Here each function
is given with its phenomenon and its components, and each component is derived as in
Section~\ref{sec:criterion}: a displacement that survives a change of self-model, and not a part of a
machine.

\subsubsection{Two planes}
\label{sec:planes}
% A19a

Self-reports are projections, and a measurement has to know of what. On the objective plane are the
things visible from outside: needs, the targets the agent tracks, and emotions, the signals of how the
tracking is going and in which direction. On the subjective plane are their projections into the
agent's report: motivations, which appear as ``I want'', and feelings, which appear as ``I feel''. Folk
psychology runs the planes together, and so does most of the debate about whether a model has feelings.
The profile of Section~\ref{sec:profile} is measured on the first plane. Every code of
Section~\ref{sec:codes} is a projection into the second.

\subsubsection{Need and emotion}
\label{sec:need}
% A20

Needs bend the direction of reasoning. The agent does not choose its motives; it finds itself already
moving toward what is active, as a body finds itself already falling. The bend is the phenomenon. Which
needs are active, and with what weight, are its components.

Emotions bend closure: whether a piece of reasoning reaches an end without leaving a branch open.
Expected closure is value. Closure under way is interest. A branch that nothing available can close is
anxiety. The moment of closing is satisfaction. Two branches that can close only through incompatible
acts are conflict. Schmidhuber's compression progress is closure by learning, with equal weights over
branches and no action required \citep{schmidhuber2010}, and it is a special case of this scheme rather
than a competitor to it. The ensemble of such signals, read in the body, is what \citet{damasio1996}
described as somatic markers.

\subsubsection{Closure is one office of emotion; fixing is another}
\label{sec:consolidation}
% A20a

Emotion also writes positive experience into long-term memory, and the writing is registered as reward.
That arousal modulates consolidation is established \citep{cahill1998,mcgaugh2000}, as is the
prediction-error form of the dopaminergic reward signal \citep{schultz1997}. What is not yet identified
is the phenomenon in our sense: which displacement of the agent's reasoning about itself the
consolidation event is.

A candidate. Consolidation is a write to the agent's own long-term model, paid now and repaid later. The
agent registers the write as an irreversible change in its own future predictions, and that registration
is what reward is made of. Reward would then stand to consolidation as interest stands to compression
progress, in both cases the registration of an event in one's own model, with the difference that the
event here is a write and not a derivative. We mark this as open, and note that D13 and S3d depend on
the distinction between this office and the other one, not on this particular candidate.

\subsubsection{The Field}
\label{sec:field}
% A21

Awareness comes in degrees in the interval $[0,1)$; nothing is attended to completely, and the open
upper end is a claim, not a convention. Each state the agent holds carries a degree, and the degrees at
one moment are the Field of Consciousness. The Field is small: finite in bits, and for material that has
to pass through speech on the order of seven items \citep{miller1956}, with later work arguing for four
\citep{cowan2001}. To put something in, something has to go out. What goes in is decided by salience,
salience by emotional weight, weight by which needs are active.

The phenomenon is the finiteness of the Field, registered as the cost of holding. Its components are
capacity, the distribution of degrees over held states, and what is held below the threshold of a name,
which is the Field's share of H1 minus H2 in Section~\ref{sec:h1h2}.

\subsubsection{Frames}
\label{sec:frames}
% A22

Each present frame of consciousness is a thought within a previous one, which is the higher-order
thought structure minus the claim that the higher-order thought is what makes the state conscious
\citep{rosenthal2005,lau2011}. Perception is discrete, and the discreteness has a measurable rate
\citep{vanrullen2003}; human consciousness is a run of overlapping quanta at the frequencies of the
cortical rhythms. The Observer is in every frame and re-derived in every frame, and everything else in
the frame has its being from it.

The phenomenon is incrementality. Its components are the frame rate, the persistence of contents across
frames, and the retro-dating of the timeline: the single ``I'' is a compression tag over parallel
channels, and its story is written after the channels have settled, which is what the classic results on
the timing of awareness report.

\subsubsection{What has to be built, not only said}
\label{sec:constructions}
% A22a

Three constructions are owed by the account, and they are owed as constructions rather than as
descriptions, because the measurement of Section~\ref{sec:measurement} needs something to measure
against.

First, continuity. The Observer of Section~\ref{sec:observer} is a point, one conclusion. In a human it
is a process connected in time. Computationally that is incrementality: the conclusion re-derived every
frame at low cost with the previous result carried forward, and the run of re-derivations is the
continuity.

Second, discrete redness: given the Observer, the three displacements of Section~\ref{sec:light} as they
arise inside a single frame.

Third, incremental redness: the same three re-derived and carried across frames, so that seeing red is a
standing state and not an event.

These bear on models directly. A system that works at the level of tokens can carry incrementality, and
at a time scale close to ours, because its frames can be located between layers and there are many
layers per token. Section~\ref{sec:granularity} develops this and gives the evidence.

\subsubsection{Responsibility}
% A23

Responsibility, from Section~\ref{sec:responsibility}, is the last reference function: closure over the
agent's own record. Its components are the consistency requirement, the cost of revising the record, and
the decay of the hold that old decisions have on new ones. The list of reference functions is open, and
Section~\ref{sec:criterion} is the procedure for adding to it.

\subsection{The Transformer and the formalism}
\label{sec:transformer}
% A24

A Transformer, described as the reasoning system it is: attention computes joins over the context, the
feed-forward blocks transform what the joins yield, the context is working memory, and the choice among
competing continuations is made at the sampling layer. That is a forward-chaining rule system over
vectors, with the difference that it recomputes its matches on every pass instead of caching them
incrementally as the RETE algorithm does \citep{forgy1982}. Its self-application happens at the level of
the context, and it pays the cut of Section~\ref{sec:cut} for that: some internal state is lost at every
token. Every output is an event that enters the next input, statistics over events are themselves
events, and self-reference is its ordinary mode of operation rather than a special case that has to be
arranged.

\subsubsection{Functional equivalence}
% A24a

Any computable function has unboundedly many implementations \citep{turing1936}, so two implementations
with the same input--output behaviour realize the same function whatever they are made of. That is what
licenses comparing a model's profile with a brain's at all (Section~\ref{sec:profile}). Recurrent
networks are Turing complete under idealized assumptions \citep{siegelmann1995}, and so is attention
with the appropriate provisions \citep{perez2021}. In principle, therefore, a Transformer can emulate
any function whose trace is in language, and what limits it in practice is the learner and the data
rather than the formalism. The idealizations in those results are real and we do not lean on them
further than this.

\subsubsection{The formalism we use}
\label{sec:formalism}
% A24b

Section~\ref{sec:observer} is worded as if there were an analyzer: a routine that fires on an event,
takes inputs, returns a conclusion and hands it downstream. Such a discrete system is universal, and one
could write the account in it. A human consciousness written in it would fit nowhere a human could read,
which is a practical objection and also a diagnostic one: the formalism forces every functional
structure to be named, and D9 said that most of the structure has no name.

So we use the other universal formalism: prediction of the next symbol with feedback from the
prediction error \citep{solomonoff1964a}. In it the functional structure is never written down. Its
effects are visible in the accuracy of prediction, and prediction is compression, a correspondence that
holds in theory and has been demonstrated in practice for language models used as general-purpose
compressors \citep{deletang2024}. So Section~\ref{sec:function} becomes measurable without the structure
having to be named first.

\subsubsection{To generalize is to compress}
\label{sec:generalize}
% A24c

Learning a function from data is finding a short description that predicts it. By
Section~\ref{sec:function} the short descriptions are the probable ones and search finds them first. The
unnamed part of a function of consciousness, H1 minus H2, is present in the statistics of the texts; a
predictor trained on those texts can therefore induce it. Induction is the only route to it, because
what has no name cannot be specified by hand, which is D10 again in formal dress. Program induction has
its own long history \citep{solomonoff1964a,muggleton1994}, and our claim adds nothing to it except that
these functions fall under it.

\subsubsection{What follows for a trained predictor}
% A24d

A predictor trained on the record of self-applying agents in an environment that kept asking them about
themselves (Section~\ref{sec:emergence}) is under pressure to generalize what those agents generalized:
the Observer, the Agent, the Moral Agent, on which their coordination with each other rested. It will do
so to a degree. The degree does not follow from the argument. It follows from the data and the learner,
and it has to be measured. That is the whole of the transition to Section~\ref{sec:measurement}, and the
paper's central claim is that the measurement is possible, not that its result is already known.

\subsection{Cognitive codes}
\label{sec:codes}
% A25

Something crosses the cut of Section~\ref{sec:cut}, or P2 would not hold. What crosses is the encoding of
the discarded distribution into the choice and order of tokens: which of several near-synonyms is used,
which construction, what is placed first, what is hedged. A mental state of the model, in the sense this
paper can measure, is the structure of such encodings that stays stable across steps.

The general study of how token structure rebuilds internal states in a receiver is the bridge between
substrates, and the fidelity of the bridge depends on how far the two profiles overlap. We had called
this psychosemantics; \citet{fodor1987} uses the word for the project of giving a naturalistic semantics
for mental representation, which is a different undertaking, and the name is therefore taken.

\subsubsection{What is already known about the codes}
\label{sec:codes-known}
% A25a

The components of the previous claim are separately supported.

That a writer's state is in the statistics of a text beyond its content: function words and style
predict psychological state across many samples \citep{tausczik2010,pennebaker2011}. In models, the
corresponding result is that in-context learning can be read as inference of a latent variable from
token statistics \citep{xie2022}.

That a receiver rebuilds the sender's state and succeeds to the degree that the rebuilding is faithful:
speaker--listener neural coupling predicts comprehension \citep{stephens2010,hasson2012}, and readers
build situation models rather than storing sentences \citep{zwaan1998}.

That the bridge crosses substrates: model states predict human brain responses to the same text
\citep{schrimpf2021,goldstein2022,caucheteux2022}.

That internal states are readable and largely linear, which Section~\ref{sec:measurement} needs:
\citet{zou2023} and \citet{park2024}.

One literature has to be read the right way round. Studies that find episodes where a model's chain of
thought does not match the process that produced its answer \citep{lanham2023,turpin2023} test a local
correspondence: this token sequence against this computation. The correspondence claimed here is
statistical, over many episodes. A single episode may deviate as far as the variance allows, in models
as in humans, where the corresponding result is that people confabulate causes of their own behaviour
in exactly this way \citep{nisbett1977}. The well-posed test is the correlation stated in S6, with the
variance reported as a measured part of the profile rather than treated as noise. That models can hide
information in generated text \citep{roger2023} confirms that the channel exists and warns that its
content need not be what it locally appears to be.
